\begin{lemma}
\label{lemma-limit-affine}
In Situation \ref{situation-descent} if $S$ is affine,
then for some $i_0 \in I$ the schemes $S_i$ for $i \geq i_0$
are affine.
\end{lemma}

\begin{proof}
By Lemma \ref{lemma-limit-quasi-affine} we may assume that $S_0$ is
quasi-affine for some $0 \in I$. Set $R_0 = \Gamma(S_0, \mathcal{O}_{S_0})$.
Then $S_0$ is a quasi-compact open of $T_0 = \Spec(R_0)$. Denote
$j_0 : S_0 \to T_0$ the corresponding quasi-compact open immersion.
For $i \geq 0$ set $\mathcal{A}_i = f_{i0, *}\mathcal{O}_{S_i}$.
Since $f_{i0}$ is affine we see that
$S_i = \underline{\Spec}_{S_0}(\mathcal{A}_i)$.
Set $T_i = \underline{\Spec}_{T_0}(j_{0, *}\mathcal{A}_i)$.
Then $T_i \to T_0$ is affine, hence $T_i$ is affine. Thus
$T_i$ is the spectrum of
$$
R_i = \Gamma(T_0, j_{0, *}\mathcal{A}_i) = \Gamma(S_0, \mathcal{A}_i) =
\Gamma(S_i, \mathcal{O}_{S_i}).
$$
Write $S = \Spec(R)$. We have $R = \colim_i R_i$
by Lemma \ref{lemma-descend-section}.
Hence also $S = \lim_i T_i$. As formation of the relative spectrum commutes
with base change, the inverse image
of the open $S_0 \subset T_0$ in $T_i$ is $S_i$.
Let $Z_0 = T_0 \setminus S_0$ and let $Z_i \subset T_i$
be the inverse image of $Z_0$. As $S_i = T_i \setminus Z_i$, it suffices
to show that $Z_i$ is empty for some $i$. Assume $Z_i$ is nonempty for all
$i$ to get a contradiction. By Lemma \ref{lemma-limit-closed-nonempty}
there exists a point $s$ of $S = \lim T_i$ which maps to a point of $Z_i$
for every $i$. But $S = \lim_i S_i$, and hence we arrive at a contradiction
by Lemma \ref{lemma-topology-limit}.
\end{proof}
